\section*{Hoofdstuk \ref{ch:b2:blood-pressure} --- Regeling van de bloeddruk en inspanning}

\begin{solution}{exo:b2:blood-pressure:1}
$\bar P - P_{\text{ven}} = Q\,R = f\,V_s\,R$. Hartfrequentie (de sinusknoop,
onder invloed van de \omterm{prop:b2:heart:control}{nervus vagus} en de sympathische zenuwen); \omterm{prop:b2:heart:cycle}{slagvolume}
(het \omterm{def:b2:heart:anatomy}{ventrikel}: zijn vulling, bepaald door de \omterm{def:b2:blood-circulation:vessels}{aders}, en zijn contractiliteit,
bepaald door de sympathische zenuwen); weerstand (het gladde spierweefsel
van de \omterm{def:b2:blood-circulation:vessels}{arteriolen}, onder invloed van de sympathische zenuwen, hormonen en
plaatselijke metabolieten).
\end{solution}

\begin{solution}{exo:b2:blood-pressure:2}
Sensoren: rekreceptoren in de sinus caroticus en de aortaboog. Centrum: de
cardiovasculaire centra van het verlengde merg. Effectoren: de \omterm{prop:b2:heart:control}{nervus vagus}
(frequentie) en de sympathische zenuwen (frequentie, contractiliteit,
arteriolaire en veneuze tonus). Negatieve terugkoppeling, ongeveer een
seconde vertraging. Zonder de sensoren blijft de gemiddelde druk over een
dag normaal, maar schommelt de druk van minuut tot minuut wild bij elke
houding en emotie.
\end{solution}

\begin{solution}{exo:b2:blood-pressure:3}
De nier (cellen van de \omterm{def:b2:blood-circulation:vessels}{arteriolen}) geeft renine af wanneer haar doorbloeding
of de aanvoer van natrium daalt of haar sympathische zenuwen vuren; renine
knipt angiotensinogeen (lever) tot angiotensine I; een enzym van de long zet
dat om in angiotensine II, dat \omterm{def:b2:blood-circulation:vessels}{arteriolen} vernauwt en de bijnierschors
aanzet aldosteron af te scheiden; aldosteron laat de nier natrium en water
vasthouden, wat het bloedvolume verhoogt.
\end{solution}

\begin{solution}{exo:b2:blood-pressure:4}
Metabole verwijding van de eigen \omterm{def:b2:blood-circulation:vessels}{arteriolen} van de spier door haar
metabolieten; een vijfvoudig hartminuutvolume, aangedreven door de
sympathische zenuwen en de spierpomp; sympathische vernauwing van de
\omterm{def:b2:blood-circulation:vessels}{arteriolen} van darm, nieren en rustende spieren, die stroom afstaan.
\end{solution}

\begin{solution}{exo:b2:blood-pressure:5}
$83\times 1.2 = \qty{100}{mmHg}$; $333\times 0.35 = \qty{117}{mmHg}$.
Weerstand gedaald met een factor $1.2/0.35 = 3.4$: straal gestegen met een
factor $3.4^{1/4} =
1.36$.
\end{solution}

\begin{solution}{exo:b2:blood-pressure:6}
$30/(1 + G)$: 10, 6 en \qty{3.3}{mmHg}. Met de halve versterking verdubbelt
de resterende daling ongeveer: ze wordt duizelig bij het opstaan, ziet grijs,
en kan flauwvallen.
\end{solution}

\begin{solution}{exo:b2:blood-pressure:7}
$\rho g h$: $1060\times 9.81\times 0.4 = \qty{4160}{Pa} =
\qty{31}{mmHg}$; $1060\times 9.81\times 1.2 = \qty{12500}{Pa} =
\qty{94}{mmHg}$. Hersenen $95 - 31 = \qty{64}{mmHg}$; enkel $95 + 94 =
\qty{189}{mmHg}$. Een kop van een giraf op \qty{2.5}{m} kost
\qty{195}{mmHg}: de giraf heeft ter hoogte van het \omterm{def:b2:heart:anatomy}{hart} een gemiddelde druk
van zo'n \qty{250}{mmHg} nodig, en dikke vaatwanden en een strakke huid in
de poten.
\end{solution}

\begin{solution}{exo:b2:blood-pressure:8}
$Q = 280/(190 - 130) = \qty{4.7}{L/min}$. Atlete: $5000/(200 - 30) =
\qty{29}{L/min}$; \omterm{prop:b2:heart:cycle}{slagvolume} $\num{29400}/190 = \qty{155}{mL}$.
\end{solution}

\begin{solution}{exo:b2:blood-pressure:9}
$35/5 = \qty{7}{mmHg}$ lager: ongeveer \qty{88}{mmHg}. Herstel van het
volume: resorptie van interstitieel vocht in de \omterm{def:b2:blood-circulation:vessels}{haarvaten} (minuten tot een
uur); minder urine onder \omterm{prop:b2:blood-pressure:raas}{antidiuretisch hormoon} en aldosteron (uren tot
dagen); dorst en drinken (uren); rode bloedcellen uit het beenmerg (weken). De
huid is bleek en koud omdat de reflex haar \omterm{def:b2:blood-circulation:vessels}{arteriolen} heeft dichtgedraaid om
de druk voor hersenen en \omterm{def:b2:heart:anatomy}{hart} te sparen.
\end{solution}

\begin{solution}{exo:b2:blood-pressure:10}
Spier bij $R = 0.5$: $95/0.5 = \qty{190}{mL/s}$, \qty{11.4}{L/min};
met de \qty{4}{L/min} van de andere organen zou het \omterm{def:b2:heart:anatomy}{hart}
\qty{15.4}{L/min} nodig hebben. Met $Q$ vast op \qty{5}{L/min} daalt de totale
weerstand van 1.14 tot $1/(1/1.14 - 1/5 + 1/0.5) =
\qty{0.37}{mmHg.s/mL}$ en de druk tot \qty{31}{mmHg} --- de hersenen zouden
uitvallen. Het lichaam verhoogt het debiet in de richting van het eerste
geval en vernauwt de andere vaatbedden, zodat de druk standhoudt en de spier
haar stroom krijgt.
\end{solution}

\begin{solution}{exo:b2:blood-pressure:11}
(a) De onvoldoende doorbloede nier las een lage druk, gaf renine af, en
angiotensine en aldosteron verhoogden de weerstand en het volume. (b) De
\omterm{def:b2:blood-pressure:baroreflex}{baroreflex} buffert veranderingen maar stelt zich opnieuw in op elk niveau
dat aanhoudt, en kan het vasthouden van volume door de nier niet
veranderen. (c) De bron van renine, en het orgaan dat een hogere druk eiste,
was weg. (d) De omzetting tot angiotensine II blokkeren: een lagere
weerstand, minder aldosteron, een lagere druk.
\end{solution}

\begin{solution}{exo:b2:blood-pressure:12}
De \omterm{def:b2:blood-pressure:baroreflex}{baroreflex} werkt in een seconde met een versterking van ongeveer vier en
een instelwaarde die naar de heersende druk verschuift: hij dempt
schommelingen en kan geen niveau vasthouden. De nier werkt over uren en
dagen, met in feite een oneindige versterking --- ze blijft uitscheiden of
vasthouden tot de druk waarop ze is ingesteld bereikt is --- en haar
instelwaarde ligt vast door haar eigen fysiologie: zij houdt het niveau vast.
Een chronische stijging betekent dus dat de instelwaarde van de nier is
verschoven, en alleen middelen die op de lus van de nier of op de weerstand
die zij eist inwerken, brengen hem omlaag.
\end{solution}

\begin{solution}{pb:b2:blood-pressure:1}
\textbf{1.} $1/R = 1/7.6 + 1/22.8 + 1/4.1 + 1/5.2 + 1/5.7 + 1/19 +
1/28.5 = 0.875$: $R = \qty{1.14}{mmHg.s/mL}$; $\bar P/Q = 95/83.3 =
1.14$.
\textbf{2.} $\bar P/R_i$: hersenen 0.75, \omterm{def:b2:heart:anatomy}{hart} 0.25, darm en lever
1.39, nieren 1.10, spieren 1.0, huid 0.30, overige \qty{0.20}{L/min};
fracties 15, 5, 28, 22, 20, 6 en \qty{4}{\%}.
\textbf{3.} $5000/70 = \qty{71}{mL}$; $5\times 50 = \qty{250}{mL}$
zuurstof per minuut.
\textbf{4.} Hersenen $0.75/1.4 = \qty{0.54}{L/min/kg}$; nieren $1.1/0.3
= \qty{3.7}{L/min/kg}$, zeven keer die van de hersenen en vijftig keer het
lichaamsgemiddelde: de nieren worden doorbloed om te filtreren, niet om te
voeden.
\textbf{5.} Met één druk die alle organen delen, trekt elk orgaan de stroom
die zijn eigen weerstand bepaalt; stromen centraal regelen zou elk orgaan
uithongeren waarvan de vraag veranderde. De druk is de gemeenschappelijke
munt.
\textbf{6.} \omterm{def:b2:heart:anatomy}{Hart}: $250\times 0.2\times 0.7 = \qty{35}{mL/min}$;
spier: $1000\times 0.2\times 0.25 = \qty{50}{mL/min}$. Het \omterm{def:b2:heart:anatomy}{hart} onttrekt al
het grootste deel van wat passeert, zodat het alleen zuurstof kan winnen
door meer stroom, en daarom verwijden zijn \omterm{def:b2:blood-circulation:vessels}{arteriolen} zodra het harder werkt.
\textbf{7.} $Q$ \qty{30}{\%} lager bij vaste $R$: $\bar P$
\qty{30}{\%} lager, \qty{28.5}{mmHg}, tot \qty{66.5}{mmHg}.
\textbf{8.} $28.5/5 = \qty{5.7}{mmHg}$: ongeveer \qty{89}{mmHg}.
\textbf{9.} $Q = 85\times 0.7\times 71 = \qty{4.2}{L/min} =
\qty{70}{mL/s}$; $R = 89.3/70 = \qty{1.27}{mmHg.s/mL}$, een stijging van
\qty{11}{\%}.
\textbf{10.} $1060\times 9.81\times 0.4 = \qty{4160}{Pa} =
\qty{31}{mmHg}$; hersenen $66.5 - 31 = \qty{35}{mmHg}$ vóór de reflex,
$89 - 31 = \qty{58}{mmHg}$ erna.
\textbf{11.} $28.5/2 = \qty{14}{mmHg}$: \qty{81}{mmHg} ter hoogte van het
\omterm{def:b2:heart:anatomy}{hart}, \qty{50}{mmHg} in de hersenen --- duizeligheid, grijs zien, bijna
flauwvallen bij elke keer opstaan uit een stoel.
\textbf{12.} Plat liggen neemt het hydrostatische verlies van
\qty{31}{mmHg} weg en laat het weggezakte \omterm{def:b2:blood-circulation:blood}{bloed} naar het \omterm{def:b2:heart:anatomy}{hart} terugvloeien;
\omterm{prop:b2:heart:cycle}{slagvolume} en druk in de hersenen herstellen binnen enkele slagen.
\textbf{13.} $45/5 = \qty{9}{mmHg}$: \qty{86}{mmHg}.
\textbf{14.} $1/R = 0.875 - 1/19 - 1/4.1 + 1/38 + 1/8.2 = 0.727$: $R =
\qty{1.38}{mmHg.s/mL}$. Hersenen $86/7.6 = \qty{11.3}{mL/s} =
\qty{0.68}{L/min}$; huid $86/38 = \qty{2.3}{mL/s} = \qty{0.14}{L/min}$;
darm $86/8.2 = \qty{10.5}{mL/s} = \qty{0.63}{L/min}$.
\textbf{15.} Met minder \omterm{def:b2:blood-circulation:blood}{bloed} en nauwere \omterm{def:b2:blood-circulation:vessels}{arteriolen} daalt de druk in de
\omterm{def:b2:blood-circulation:vessels}{haarvaten} over het hele \omterm{def:b2:blood-circulation:vessels}{haarvat} onder de \omterm{thm:b2:blood-circulation:oncotic}{oncotische druk}, zodat vocht wordt
geresorbeerd in plaats van gefiltreerd; het \omterm{def:b2:blood-circulation:blood}{plasma} wordt verdund en de
\omterm{def:b2:blood-circulation:blood}{hematocriet} daalt.
\textbf{16.} Aldosteron, via renine en angiotensine, in gang gezet door de
lage doorbloeding van de nier en de sympathische prikkeling; \omterm{prop:b2:blood-pressure:raas}{antidiuretisch
hormoon}, in gang gezet door de daling van het volume en de stijging van de
plasmaconcentratie.
\textbf{17.} $5\times 10^{12}$ cellen bij $2.5\times 10^{6}$ per seconde:
$2\times 10^{6}$ s, ongeveer drie weken.
\textbf{18.} Onder ongeveer \qty{60}{mmHg} is de reflexcurve vlak: de vaten
zijn zo nauw en het \omterm{def:b2:heart:anatomy}{hart} zo snel als ze kunnen zijn, de versterking is nul,
en elk verder verlies verlaagt de druk met zijn volle waarde; de onvoldoende
doorbloede weefsels geven dan zuur af en verwijden, en de druk stort in.
\textbf{19.} Spieren $110/0.33 = \qty{333}{mL/s} = \qty{20}{L/min}$;
huid $110/5 = \qty{22}{mL/s} = \qty{1.3}{L/min}$.
\textbf{20.} Weerstand van de hersenen $8.8$ (stroom 0.75 bij 110), \omterm{def:b2:heart:anatomy}{hart} $6.6$
(stroom 1.0): $1/R = 1/8.8 + 1/6.6 + 1/8.2 + 1/10.4 + 1/0.33 + 1/5 +
1/28.5 = 3.75$: $R = \qty{0.27}{mmHg.s/mL}$; $Q = 110/0.27 =
\qty{410}{mL/s} = \qty{25}{L/min}$.
\textbf{21.} $\num{24700}/180 = \qty{137}{mL}$.
\textbf{22.} $24.7\times(200 - 40) = \qty{3950}{mL/min}$, zestien keer de
\qty{250}{mL/min} in rust.
\textbf{23.} Stroom $\times 4.9$, extractie $160/50 = \times 3.2$;
$4.9\times 3.2 = 16$.
\textbf{24.} Het bevel van de motorische schors en de signalen van de
receptoren van de spieren herwegen de input van de \omterm{def:b2:blood-pressure:baroreflex}{baroreceptoren} in het
verlengde merg, zodat de reflex \qty{110}{mmHg} verdedigt in plaats van 95
--- de instelwaarde wordt verhoogd, de lus niet uitgeschakeld.
\textbf{25.} Resterende daling bij het opstaan \qty{5.7}{mmHg}; \qty{86}{mmHg}
na de bloeding; \qty{25}{L/min} en \qty{4}{L} zuurstof per minuut bij
inspanning.
\end{solution}
